\subsection{Cas limites}
\label{sec:cas_limites}

\begin{enumerate}
    \item \textbf{Succession CTO intégralement exonérée.} Si \(B_{CTO} = 0\), alors \(\Delta H > 0\). Le CTO ne peut être désavantagé tant qu'il n'est pas imposé.
    \item \textbf{Succession AV intégralement exonérée.} Si \(B_{AV} = 0\), alors
    \begin{equation}
        \Delta H = (C_{CTO} - C_{AV}) + PS - s f(B_{CTO}).
    \end{equation}
    L'AV ne devient plus avantageuse que si \(s f(B_{CTO}) \geq (C_{CTO} - C_{AV}) + PS\).
\end{enumerate}

Dans la pratique française, ce second cas est particulièrement informatif: pour de nombreux ménages transmettant en ligne directe, l'abattement propre à l'assurance-vie réduit souvent la base imposable de l'AV, tandis que le CTO demeure exposé aux droits sur une base positive, mais, l'accumulation du capital est plus importante dans le cas du CTO et les plus values sont exonérés, il s'agit du cas d'intérêt que nous cherchons à étudier dans cette étude pour déterminer le signe de \(\Delta H\).
